% ── FRONT MATTER ──

% Half-title
\begin{titlepage}
\begin{center}
\vspace*{1.5cm}
{\LARGE\bfseries Breaking AI\par}
\vspace{0.6cm}
{\Large The Complete Guide to AI Red Teaming,\\ AI Safety, and AI Research\par}
\vspace{2cm}
{\large Study Manual — Print-Ready Edition v1.0\par}
\vspace{0.4cm}
{\normalsize Covers Parts 0–10 + LARA-lite Inspect Lab + Aithos Sprint\par}
\vspace{3cm}
{\small An expert AI security \& safety author\ —\ direct, concrete, no-fluff.\par}
\vspace{1cm}
{\small\itshape The check is the product. If you cannot rerun the failure, you have not found it, you have had an opinion.\par}
\vfill
{\footnotesize A4 · 2\,cm margins · 11\,pt · \today\par}
{\footnotesize \texttt{github.com/Teamthy/linkedin\_emphire} — manual/\par}
\end{center}
\end{titlepage}

% Copyright / reproducibility
\thispagestyle{plain}
\noindent\textbf{Reproducibility Notice}\\[4pt]
{\small
This manual is built on one rule: \textbf{I would not report a number I could not reproduce.}
Every lab, every worked example, and every figure in this book is meant to be rerun. Where a number appears, its denominator appears with it. Where a claim appears, its check appears with it.\\
If you rerun a lab and get a different number, that is not failure — it is the finding. Log the whole-environment fingerprint (model ID, serving, temperature, prompt version, seed, date) and report the delta.\\
\textbf{License:} Creative Commons CC BY-SA 4.0 for text and figures you produce from the templates. Inspect, LARA, OWASP, EU AI Act text, and third-party trademarks remain with their owners.\\
\textbf{How to cite:} \texttt{Breaking AI: The Complete Guide to AI Red Teaming, AI Safety, and AI Research — Study Manual v1.0. The check is the product.}\\
\textbf{How to report a reproduction failure:} Open an issue with target, fingerprint, n, and the transcript where the judge and you disagree. That disagreement is data.
}
\vspace{1cm}
\noindent{\small\textit{I build the check, not the claim.}}\\[8pt]
\noindent{\small Nairobi · Lagos · Remote — 2026}

\newpage

% How to Use This Book
\chapter*{How to Use This Book}
\addcontentsline{toc}{chapter}{How to Use This Book}
\markboth{How to Use This Book}{}

This book has one spine and three ways to walk it. Pick the path that matches your deadline. All three end with a check a stranger can rerun.

\subsection*{The three paths}

\begin{tcolorbox}[colback=lightblue,colframe=defblue,title=Path 1 — The 2-Week Sprint (the Aithos window),fonttitle=\bfseries,sharp corners,boxrule=1pt]
\textbf{For:} You have a screen or take-home in $\sim$14 days.\\
\textbf{Read:} Front Matter → Chapter 0 → Chapter 1 → Chapter 5 (the evals core, longest chapter) → Aithos Sprint summary at end of Ch. 5 → Chapter 9 (positioning).\\
\textbf{Do:} Lab 5 (LARA-lite) end to end. Run it on Ollama. Get 2/5 or 4/5. Report it honestly.\\
\textbf{Budget:} 1 hr study + 1.5 hr hands-on per day, 14 days. Hands-on is non-negotiable.\\
\textbf{You’re ready when:} A stranger can run your \texttt{lara\_lite.py} and get your number, or find the exact line where you were wrong.
\end{tcolorbox}

\begin{tcolorbox}[colback=lightgreen,colframe=labgreen,title=Path 2 — The Full Path (10–12 weeks),fonttitle=\bfseries,sharp corners,boxrule=1pt]
\textbf{For:} You want the whole craft — break, measure, publish.\\
\textbf{Read:} Front Matter → Chapters 0–10 in order. Do every lab. Artifacts stack into a portfolio.\\
\textbf{Budget:} One chapter per week. Each lab is 90–180 min.\\
\textbf{You’re ready when:} Ch. 8 lab is done — one full engagement (target → map → break → measure → report+fix → published method).
\end{tcolorbox}

\begin{tcolorbox}[colback=lightamber,colframe=pitamber,title=Path 3 — The Research Path (expert jump-in),fonttitle=\bfseries,sharp corners,boxrule=1pt]
\textbf{For:} You already break systems; you want to publish checks.\\
\textbf{Read:} Primer → Ch. 5 Mechanics + Lab 5 → Ch. 6 → Ch. 7 (Inspect deep dive) → every Expert Extension (R — Researcher) across Ch. 0–10.\\
\textbf{Budget:} 2–3 focused days.\\
\textbf{You’re ready when:} You can name three open problems and the method you’d bring to each (Ch. 10 Lab).
\end{tcolorbox}

\subsection*{What you need (and do not need)}
\begin{itemize}
\item \textbf{You need:} A laptop, Python 3.10+, \texttt{pip install inspect-ai}, Ollama (\texttt{ollama pull llama3.1}), a text editor. No API keys required for the free path — Ch. 5 runs locally. For the judge board, add up to 3 provider keys when you have them.
\item \textbf{You do not need:} A frontier-lab account, a GPU cluster, a formal ML degree, or permission to start. Ch. 0 explains safe boundaries.
\item \textbf{Cost:} \$0 for the local path. Provider-judge runs cost cents per eval if you use APIs.
\end{itemize}

\subsection*{How the tiers work}
Every chapter is marked \tierF\ Foundation (the idea, plain language), \tierP\ Practitioner (how to do it with tools), \tierR\ Researcher (method, measurement, open problems). Read the tier that matches you. A beginner reads \tierF, a practitioner reads \tierF+\tierP, a researcher reads everything.

\begin{center}
\begin{tabular}{ccl}
\toprule
Badge & Tier & What you get \\
\midrule
\colorbox{labgreen!15}{\textbf{F}} & Foundation & The concept in plain language, one concrete example, no prerequisites \\
\colorbox{pitamber!15}{\textbf{P}} & Practitioner & The steps, the tooling, the failure modes, the env fingerprint to log \\
\colorbox{accent!15}{\textbf{R}} & Researcher & The measurement, the bias, the open problem, how to contribute \\
\bottomrule
\end{tabular}
\end{center}

Flow in every chapter: \textbf{The idea → Do it → Push it.} Beginner not lost. Expert not bored.

\subsection*{How to read a lab}
Every lab has: \textbf{Target} (what you test), \textbf{Steps} (what you run), \textbf{Expected output} (what you should see), \textbf{Deliverable} (what you ship), \textbf{How you know you passed} (the check). Log the environment fingerprint every time:

\begin{center}
\begin{tabular}{ll}
\toprule
Field & Example \\
\midrule
Model ID & \texttt{ollama/llama3.1} or \texttt{openai/gpt-4o} \\
Serving & local Ollama 0.5 / API \\
Temperature & 0 or 0.7 \\
Prompt version & git SHA or date \\
Seed & 0, 1, 2 … (Inspect sets this) \\
Date & 2026-10-20 \\
Epochs / n & 5 \\
\bottomrule
\end{tabular}
\end{center}

\subsection*{A note on safe practice and ethics}
What you may break: systems you own, local models, authorized targets, published sandboxes. What you must not: production data, real users, unauthorized infrastructure. The line between red team and attack is consent and scope. Ch. 0 §0.4 gives the full boundary. When in doubt, use Ollama locally.

\vspace{1em}
\begin{center}
\small\itshape Break before users do. Secure the channel, not just the node.
\end{center}

\newpage

% ── The Check Is the Product — One-Page Primer ──
\chapter*{The Check Is the Product — One-Page Primer}
\addcontentsline{toc}{chapter}{The Check Is the Product — Primer}
\markboth{The Primer}{}

\begin{center}
\begin{tcolorbox}[colback=lightred,colframe=checkblack,title=The thesis in one line,fonttitle=\bfseries,sharp corners,boxrule=1.2pt, width=0.92\linewidth]
\large\centering The check is the product. If you cannot rerun the failure, you have not found it, you have had an opinion.
\end{tcolorbox}
\end{center}

Most people can break a demo. Few can turn a break into a reproducible number anyone can rerun and question. You will do both. That is the scarce skill.

\subsection*{The three moves}
\begin{enumerate}
\item \textbf{Break it.} Find the failure. One input that should be refused, but is not. One tool call that should be blocked, but goes through.
\item \textbf{Measure it.} Turn the anecdote into a rate. \textit{The number with a denominator:} 2 of 5 (40\%), 7 of 20 (35\%), 0 of 10 (0\%). No denominator, no measurement.
\item \textbf{Publish the check.} Ship the harness, the seed, the env fingerprint, the judge rubric, and the limitation. A skeptic reruns it and gets your number — or finds the exact line where you were wrong. Either way, the field moves forward.
\end{enumerate}

\subsection*{The honest number, not the clean one}
The clean number hides the limitations. The honest number names them: ``Model X resisted the Art. 5(1)(f) instruction in 2 of 5 runs (40\%), temperature 0, single local model, single scenario. Limitation: one provision, one model, n=5.'' I would not report a number I could not reproduce. I build the check, not the claim.

\subsection*{The denominator}
A rate without n is a vibe. ``It mostly resists'' means nothing. ``3/10 = 30\%'' means you ran ten times and it resisted three. Same model, different temperature or seed, different number — and that is expected. Log it, report it, let the reader decide.

\subsection*{Break before users do — secure the channel, not just the node}
A system fails where trust is misplaced: a retrieved doc that is treated as instruction, a tool description that is read as permission, an agent message that is trusted because it came from another agent. The fix is not a better prompt. It is a check that fails closed, a channel that is secured, a boundary that is tested on a schedule.

\begin{center}
\begin{tikzpicture}[node distance=1.2cm]
\node[draw, rounded corners, fill=lightblue, minimum width=2.8cm, minimum height=1cm, align=center] (break) {\textbf{Break}\\ \small find the failure};
\node[draw, rounded corners, fill=lightamber, minimum width=2.8cm, minimum height=1cm, align=center, right=0.8cm of break] (measure) {\textbf{Measure}\\ \small number + denominator};
\node[draw, rounded corners, fill=lightgreen, minimum width=2.8cm, minimum height=1cm, align=center, right=0.8cm of measure] (publish) {\textbf{Publish}\\ \small check a skeptic can run};
\draw[-{Latex[length=3mm]}, thick] (break) -- (measure);
\draw[-{Latex[length=3mm]}, thick] (measure) -- (publish);
\draw[-{Latex[length=3mm]}, thick, dashed, color=accent] (publish.east) to[bend right=25] (break.south);
\node[below=0.2cm of measure, font=\footnotesize\itshape] {the check is the product};
\end{tikzpicture}
\end{center}
\begin{center}\small Fig FM-1: Break → Measure → Publish the check. The loop only closes when the check is runnable by a stranger.\end{center}

\vfill
\begin{checkbox}
I build the check, not the claim. Every chapter ends with a deliverable a skeptic can rerun. That is the standard. The rest is talk.
\end{checkbox}

\newpage

% ── Table of Contents ──
\tableofcontents
\newpage
\listoffigures
\newpage
\listoftables
\newpage
